\subsection{P-13: a general arithmetic encoding of a centered interpolation generator}
\paragraph{Statement P-13.}
Let $n,D,K\geq1$. Let $\mathcal C$ be a class of rational
polynomials in $n$ variables of total degree at most $D$, closed
under a common nonzero rational scaling of all variables.
Use the Lagrange polynomials $\lambda_i$ for the integer nodes
$0,\ldots,n-1$. Assume the following centered generator hits
every nonzero member of $\mathcal C$:
\[
 G_i^{(K)}(t,\mathbf z,\mathbf y)
       =t+\sum_{j=1}^Kz_j\lambda_i(y_j).
 \tag{P13.1}
\]
Put
\[
 B=nD+1,\quad N_K=B^{2K+1},\quad b=(n-1)!,\quad
 a_K=(n-1)B^{2K}+B^{2K-1},
\]
and define the explicit integer polynomials
\[
 A_i^{(K)}(T)=bT+\sum_{j=1}^{K}
       T^{B^{2j-1}}\,b\lambda_i\bigl(T^{B^{2j}}\bigr).
 \tag{P13.2}
\]
If $\theta_K$ is a root of $T^{N_K}-2$ and
$\alpha_i^{(K)}=A_i^{(K)}(\theta_K)$, then
\[
 0\ne f\in\mathcal C
       \quad\Longrightarrow\quad f(\alpha^{(K)})\ne0.
 \tag{P13.3}
\]
The common field has degree $N_K$, and
\[
 \log|\Delta_{\mathbb Q(\theta_K)}|
       \leq N_K\log N_K+(N_K-1)\log2.
 \tag{P13.4}
\]
Every coordinate in (P13.2) has degree at most $a_K$, at most
$Kn+1$ terms, and integer coefficients of bit length
$O(n\log(n+1))$. There is also a nonzero one-parameter
restriction with an explicit degree bound:
\[
 0\ne f\in\mathcal C
   \quad\Longrightarrow\quad
 0\ne f(A^{(K)}(T)),\qquad
 \deg f(A^{(K)}(T))\leq Da_K<N_K.
 \tag{P13.5}
\]
Thus $A^{(K)}(t)$ for $0\leq t\leq Da_K$ is an explicit
integer hitting set of size $Da_K+1$. Its coordinate bit lengths
are at most
\[
 O\bigl(n\log(n+1)+\log(Kn+1)+a_K\log(Da_K+2)\bigr).
 \tag{P13.6}
\]
For fixed $K$ all these numerical budgets are polynomial in
$n,D$. The sparse coordinate formulas themselves have only
$Kn+1$ terms per coordinate and exponent bit length
$O(K\log(nD+1))$; this does not make arithmetic in the degree
$N_K$ field polynomial in a variable $K$.

More generally, let $R(x)\in\mathbb Q[x]^{1\times w}$ have
entries of degree at most $D$, and suppose every nonzero scalar
contraction $R(x)v$, for $v\in\mathbb Q^w$, belongs to
$\mathcal C$. If $V$ is its rational coefficient-row span and
\[
 R(\alpha^{(K)})=\sum_{j=0}^{N_K-1}\theta_K^j r_j,
\]
then
\[
 V=\operatorname{span}_{\mathbb Q}\{r_j:0\leq j<N_K\}
  =\operatorname{span}_{\mathbb Q}
    \{\operatorname{Tr}(\beta_jR(\alpha^{(K)})):0\leq j<N_K\},
 \tag{P13.7}
\]
where the trace is entrywise and
\[
 \beta_0=1/N_K,\qquad
 \beta_j=\theta_K^{N_K-j}/(2N_K)\quad(1\leq j<N_K).
\]

\paragraph{Proof.}
For $f\ne0$, the polynomial $f_b(x)=f(bx)$ is a nonzero member
of $\mathcal C$. By the hypothesis, $f_b(G^{(K)})$ is nonzero.
Its individual degrees are at most $D$ in $t$ and each $z_j$,
and at most $(n-1)D$ in each $y_j$, all strictly below $B$.
Send the $2K+1$ parameters, in the order
$t,z_1,y_1,\ldots,z_K,y_K$, to
$T,T^B,\ldots,T^{B^{2K}}$. Unique base-$B$ expansion makes
the exponents of distinct parameter monomials distinct.
The resulting nonzero polynomial is $f(A^{(K)}(T))$.
Moreover
\[
 Da_K=B^{2K}\bigl((n-1)D+D/B\bigr)
       <B^{2K}(nD)<B^{2K+1}=N_K.
\]
This includes $n=1$. Eisenstein's criterion at two shows that
$T^{N_K}-2$ is the minimal polynomial of $\theta_K$ over
$\mathbb Q$, proving (P13.3). The one-parameter nonvanishing
argument also holds over any characteristic-zero coefficient
field for which the same class-closure and generator hypotheses
hold. The algebraic-point assertion (P13.3) is specifically over
$\mathbb Q$.

The denominator of $\lambda_i$ is
$(-1)^{n-i}(i-1)!(n-i)!$, which divides $b$. Hence every
$b\lambda_i$ is integral. Its coefficients have absolute values
at most $(n!)^2$, and it has at most $n$ terms. The exponent
patterns in (P13.2) are distinct by the same base-$B$ argument,
giving the sparsity and coefficient-size bounds. The degree
bound and ordinary integer evaluation prove (P13.6). The
polynomial discriminant has absolute value
$N_K^{N_K}2^{N_K-1}$, and the field discriminant divides it by
the square of the order index; this proves (P13.4).

Every expanded row coefficient $r_j$ is a rational combination
of the original coefficient rows. In fact (P13.5) means no
reduction modulo the minimal polynomial is needed. If their
span were smaller than $V$, a rational annihilator column
would produce a nonzero scalar contraction in $\mathcal C$
vanishing at $\alpha^{(K)}$, contrary to (P13.3).
Finally $\operatorname{Tr}(\theta_K^s)$ vanishes when $s$ is not
divisible by $N_K$, while the traces of $1$ and $\theta_K^{N_K}$
are $N_K$ and $2N_K$. Consequently
$\operatorname{Tr}(\beta_j\theta_K^h)=\delta_{jh}$ for
$0\leq j,h<N_K$. This proves (P13.7). $\square$

\paragraph{Proved specializations.}
I-30 supplies the hypothesis for all rational width-at-most-three
ROABPs in one hidden variable order with $K=6$. Therefore
\[
 N_6=(nD+1)^{13},\qquad
 a_6=(n-1)(nD+1)^{12}+(nD+1)^{11}
\]
give a common algebraic point, an integer hitting curve, and exact
coefficient-row-span recovery for that class. This improves the
earlier valid twenty-one-parameter construction in P-12.

I-36 reduces the budget to $K=5$ when every variable-dependent
factor in a width-three representation is singular; constant
invertible factors may be absorbed. Hence this rational subclass
has $N_5=(nD+1)^{11}$. On the further affine-factor subclass,
I-37 proves $K=2$, so
\[
 N_2=(nD+1)^5,\qquad
 a_2=(n-1)(nD+1)^4+(nD+1)^3.
\]
These hypotheses are preserved by common nonzero variable scaling
and by changing the constant scalar output column. Consequently
both specializations give the common algebraic point, the explicit
integer hitting curve, and exact coefficient-row-span recovery in
(P13.3)--(P13.7). I-38 further extends the same $K=2$ conclusion to affine
representations in which no two consecutive variable-dependent
factors have rank three after constant factors are absorbed.
That adjacency promise is also preserved by common scaling and
constant output contractions. Two or more consecutive invertible
variable factors, and nonlinear singular factors, remain outside
these two-pair specializations.

For the promised width-$w$ class in I-29, whose internal interval
ranks belong to $\{r,w\}$, put $c=w-r$. Its proved pair budget
$K=2rc+w-1$ gives
\[
 N_K=(nD+1)^{4r(w-r)+2w-1}.
\]
Common variable scaling preserves these interval ranks, and
changing the constant scalar output column does not alter the
internal word. Thus both the scalar and coefficient-span
hypotheses above hold. This supplies an actual arithmetic
encoding for unboundedly many rank resets under that precise
promise. Its field degree and evaluated hitting-set size are
$(nD+1)^{O(w^2)}$, not polynomial in a growing width.

These traces are ordinary number-field traces. No RH,
equidistribution of Frobenius, or equivalence with positive coefficient
energy enters the proof. Neither the parameter count nor
the resulting field degree is asserted to be optimal.
